\documentclass[]{article}
\usepackage{lmodern}
\usepackage{amssymb,amsmath}
\usepackage{ifxetex,ifluatex}
\usepackage{fixltx2e} % provides \textsubscript
\ifnum 0\ifxetex 1\fi\ifluatex 1\fi=0 % if pdftex
  \usepackage[T1]{fontenc}
  \usepackage[utf8]{inputenc}
\else % if luatex or xelatex
  \ifxetex
    \usepackage{mathspec}
  \else
    \usepackage{fontspec}
  \fi
  \defaultfontfeatures{Ligatures=TeX,Scale=MatchLowercase}
    \setsansfont[]{Calibri Light}
\fi
% use upquote if available, for straight quotes in verbatim environments
\IfFileExists{upquote.sty}{\usepackage{upquote}}{}
% use microtype if available
\IfFileExists{microtype.sty}{%
\usepackage{microtype}
\UseMicrotypeSet[protrusion]{basicmath} % disable protrusion for tt fonts
}{}
\usepackage[margin=0.5cm]{geometry}
\usepackage{hyperref}
\hypersetup{unicode=true,
            pdftitle={Rapport d'itinéraire},
            pdfauthor={Pivot Science},
            pdfborder={0 0 0},
            breaklinks=true}
\urlstyle{same}  % don't use monospace font for urls
\usepackage{longtable,booktabs}
\usepackage{graphicx,grffile}
\makeatletter
\def\maxwidth{\ifdim\Gin@nat@width>\linewidth\linewidth\else\Gin@nat@width\fi}
\def\maxheight{\ifdim\Gin@nat@height>\textheight\textheight\else\Gin@nat@height\fi}
\makeatother
% Scale images if necessary, so that they will not overflow the page
% margins by default, and it is still possible to overwrite the defaults
% using explicit options in \includegraphics[width, height, ...]{}
\setkeys{Gin}{width=\maxwidth,height=\maxheight,keepaspectratio}
\IfFileExists{parskip.sty}{%
\usepackage{parskip}
}{% else
\setlength{\parindent}{0pt}
\setlength{\parskip}{6pt plus 2pt minus 1pt}
}
\setlength{\emergencystretch}{3em}  % prevent overfull lines
\providecommand{\tightlist}{%
  \setlength{\itemsep}{0pt}\setlength{\parskip}{0pt}}
\setcounter{secnumdepth}{0}
% Redefines (sub)paragraphs to behave more like sections
\ifx\paragraph\undefined\else
\let\oldparagraph\paragraph
\renewcommand{\paragraph}[1]{\oldparagraph{#1}\mbox{}}
\fi
\ifx\subparagraph\undefined\else
\let\oldsubparagraph\subparagraph
\renewcommand{\subparagraph}[1]{\oldsubparagraph{#1}\mbox{}}
\fi

%%% Use protect on footnotes to avoid problems with footnotes in titles
\let\rmarkdownfootnote\footnote%
\def\footnote{\protect\rmarkdownfootnote}

%%% Change title format to be more compact
\usepackage{titling}

% Create subtitle command for use in maketitle
\providecommand{\subtitle}[1]{
  \posttitle{
    \begin{center}\large#1\end{center}
    }
}

\setlength{\droptitle}{-2em}

  \title{Rapport d'itinéraire}
    \pretitle{\vspace{\droptitle}\centering\huge}
  \posttitle{\par}
    \author{Pivot Science}
    \preauthor{\centering\large\emph}
  \postauthor{\par}
      \predate{\centering\large\emph}
  \postdate{\par}
    \date{04 septembre 2021}


\begin{document}
\maketitle

\fontsize{9}{9}
\selectfont

\hypertarget{le-trajet-du-parcours}{%
\subsection{Le trajet du parcours}\label{le-trajet-du-parcours}}

\begin{center}\includegraphics{reportCourse_files/figure-latex/tmp.data$road.fig-1} \end{center}

\begin{longtable}[]{@{}llll@{}}
\toprule
\begin{minipage}[b]{0.37\columnwidth}\raggedright
Localisation\strut
\end{minipage} & \begin{minipage}[b]{0.11\columnwidth}\raggedright
Distance\strut
\end{minipage} & \begin{minipage}[b]{0.25\columnwidth}\raggedright
Duration\strut
\end{minipage} & \begin{minipage}[b]{0.15\columnwidth}\raggedright
Altitude\strut
\end{minipage}\tabularnewline
\midrule
\endhead
\begin{minipage}[t]{0.37\columnwidth}\raggedright
Point de départ :
\textbf{47.4690034792187},\textbf{-21.2781376526266}\strut
\end{minipage} & \begin{minipage}[t]{0.11\columnwidth}\raggedright
\textbf{0 km}\strut
\end{minipage} & \begin{minipage}[t]{0.25\columnwidth}\raggedright
Min : \textbf{50 min} (3.26 km/h)\strut
\end{minipage} & \begin{minipage}[t]{0.15\columnwidth}\raggedright
Min : \textbf{590 m}\strut
\end{minipage}\tabularnewline
\begin{minipage}[t]{0.37\columnwidth}\raggedright
Point d'arrivée :
\textbf{47.4849336089781},\textbf{-21.2640605714461}\strut
\end{minipage} & \begin{minipage}[t]{0.11\columnwidth}\raggedright
\textbf{2.96 km}\strut
\end{minipage} & \begin{minipage}[t]{0.25\columnwidth}\raggedright
Max : \textbf{54 min} (2.98 km/h)\strut
\end{minipage} & \begin{minipage}[t]{0.15\columnwidth}\raggedright
Max : \textbf{720 m}\strut
\end{minipage}\tabularnewline
\bottomrule
\end{longtable}

\hypertarget{information-sur-le-trajet}{%
\subsection{Information sur le trajet}\label{information-sur-le-trajet}}

\begin{itemize}
\item
  Rivière : \textbf{0.09 km} \textbar{} à \textbf{0.98 km} il y a une
  rivière à \textbf{vaste étendue}
\item
  Forêt dense : \textbf{0.29 km}
\item
  Végétation herbeuse : \textbf{2.17 km}
\item
  Rizière : \textbf{0.1 km}
\item
  Zone d'habitation : \textbf{0.1 km}
\end{itemize}


\end{document}
